\UseRawInputEncoding
\documentclass[11pt,a4paper]{article}
\usepackage[utf8]{inputenc}
\usepackage[T1]{fontenc}
\usepackage[portuguese]{babel}
\usepackage{graphicx}
\usepackage{enumitem}
\usepackage{multicol}
\usepackage{booktabs}
\usepackage{xcolor}
\usepackage{titlesec}
\usepackage{fancyhdr}
\usepackage{geometry}
\usepackage{caption}
\usepackage{microtype} % Melhora a justificação e espaçamento entre caracteres
\usepackage[colorlinks=true, linkcolor=primaryColor, urlcolor=secondaryColor]{hyperref}
\usepackage{underscore}

% --- Definição de Cores Profissionais ---
\definecolor{primaryColor}{RGB}{26, 54, 93}     % Azul Escuro / Naval
\definecolor{secondaryColor}{RGB}{43, 108, 176} % Azul Intermédio
\definecolor{textColor}{RGB}{45, 55, 72}        % Cinza Escuro para texto (menos agressivo que preto puro)
\definecolor{lightGray}{RGB}{247, 250, 252}    % Fundo para caixas ou tabelas

% --- Configuração das Margens (Otimização de Espaço) ---
\geometry{
    a4paper,
    top=30mm,
    bottom=30mm,
    left=25mm,
    right=25mm,
    headheight=15pt,
    footskip=15mm
}

% --- Customização do Texto Global ---
\color{textColor}
\renewcommand{\baselinestretch}{1.15} % Espaçamento entre linhas ligeiramente superior

% --- Configuração de Cabeçalhos e Rodapés ---
\pagestyle{fancy}
\fancyhf{}
\fancyhead[L]{\small\color{gray}Relatório de Análise de Base de Dados: Cineteatro}
\fancyhead[R]{\small\color{gray}Grupo P2G9}
\fancyfoot[C]{\small\color{textColor}\thepage}
\renewcommand{\headrulewidth}{0.5pt}
\renewcommand{\headrule}{\hbox to\headwidth{\color{gray!50}\leaders\hrule height \headrulewidth\hfill}}

% --- Estilização das Secções ---
\titleformat{\section}
  {\Large\bfseries\color{primaryColor}}
  {\thesection}{1em}{}[{\color{secondaryColor!40}\titlerule[1.5pt]}]

\titleformat{\subsection}
  {\large\bfseries\color{secondaryColor}}
  {\thesubsection}{1em}{}

% --- Configuração de Legendas (Captions) ---
\captionsetup{
    labelfont={bf,color=primaryColor},
    font=small,
    labelsep=quad,
    margin=10pt
}

% --- Ajuste nos Espaçamentos das Listas ---
\setlist[itemize]{leftmargin=*, topsep=4pt, itemsep=3pt, parsep=0pt, partopsep=0pt}

\graphicspath{{Images/}}

\usepackage{listings}
\usepackage{xcolor}

\definecolor{codegray}{rgb}{0.95,0.95,0.95}

\lstset{
    backgroundcolor=\color{codegray},
    basicstyle=\ttfamily\small,
    frame=single,
    breaklines=true,
    keywordstyle=\color{blue},
    commentstyle=\color{green!50!black},
    stringstyle=\color{red},
    showstringspaces=false,
    tabsize=4
}

\begin{document}

% ==============================================================================
% PÁGINA DE ROSTO CUSTOMIZADA (DESIGN ELEGANTE E MODERNO)
% ==============================================================================
\begin{titlepage}
    \centering
    
    \vspace*{1cm}

    {\color{secondaryColor}\rule{\linewidth}{2pt}\par}
    \vspace{0.6cm}

    {\huge \bfseries \color{primaryColor}
    Relatório de Análise de Sistema\\
    de Base de Dados: Cineteatro\par}

    \vspace{0.4cm}

    {\large \color{secondaryColor}
    Gestão Integrada de Eventos, Bilheteira, Bar e Recursos Humanos\par}

    \vspace{0.2cm}

    {\color{secondaryColor}\rule{\linewidth}{2pt}\par}

    \vspace{2cm}

    {\Large \textsc{Unidade Curricular de BD}\par}

    \vfill

    \begin{tabular}{rll}
        \textbf{Grupo P2G9} & Ricardo Carmo   & (112657) \\
                            & Daniel Oliveira & (115359)
    \end{tabular}

    \vfill

    {\large \color{textColor!70} \today\par}

    \vspace*{1cm}

\end{titlepage}

\clearpage
\tableofcontents
\clearpage

% ==============================================================================
% CONTEÚDO DO RELATÓRIO
% ==============================================================================

\section{Resumo Executivo}
Este relatório resume a conceção e implementação de um sistema integrado para a gestão de um Cineteatro. O sistema inclui gestão de eventos, bilheteira, bar e recursos humanos, suportado por um modelo relacional normalizado e por uma aplicação cliente em WinForms que comunica com o SQL Server através de procedimentos armazenados.

\section{Introdução}
O trabalho foi realizado no âmbito da unidade curricular de Base de Dados. O objetivo principal foi conceber uma solução fiável e mantível que assegure integridade dos dados, permita operações transaccionadas e ofereça uma interface operacional para utilizadores não técnicos (bilheteira e bar).

\section{Análise de Requisitos}

\subsection{Requisitos Funcionais}
Os requisitos funcionais detalham os comportamentos, as ações e as operações essenciais que a base de dados deve suportar para satisfazer as necessidades de negócio do Cineteatro:
\begin{itemize}
    \item \textbf{Gestão de Pessoas:} Registo unificado de clientes e funcionários com dados pessoais partilhados (NIF, nome, contacto, morada). A herança de dados garante que não haja duplicações de informação básica.
    \item \textbf{Eventos:} Criação e gestão detalhada de espetáculos, incluindo a sua duração, data e hora, tipo de evento e o local correspondente. Cada evento é modelado de forma a suportar a associação de múltiplos artistas.
    \item \textbf{Bilheteira:} Emissão, controlo e venda de bilhetes de acesso para eventos específicos. Cada bilhete é processado e emitido por um funcionário específico da bilheteira.
    \item \textbf{Bar:} Venda de produtos consumíveis no bar do cineteatro, mantendo um registo detalhado de compras indexado por cliente e associado ao funcionário responsável pela transação.
    \item \textbf{Recursos Humanos:} Atribuição e controlo de funcionários nas respetivas secções, tais como bares específicos e eventos de organização, permitindo gerir as equipas de forma dinâmica.
\end{itemize}

\subsection{Requisitos Não-Funcionais}
Os requisitos não-funcionais especificam as qualidades técnicas, restrições, critérios de desempenho e integridade que asseguram o correto funcionamento do sistema:
\begin{itemize}
    \item \textbf{Integridade:} Utilização rigorosa de chaves estrangeiras (\textit{Foreign Keys}) para garantir a integridade referencial e a consistência absoluta dos dados entre todas as entidades relacionadas.
    \item \textbf{Escalabilidade:} Aplicação da diretiva \texttt{IDENTITY} para o auto-incremento de IDs únicos, assegurando que o sistema possa escalar em volume de registos sem conflitos de chaves.
    \item \textbf{Normalização:} Estruturação de dados em conformidade com a Terceira Forma Normal (3NF), eliminando redundâncias prejudiciais e evitando anomalias de inserção, atualização ou eliminação.
    \item \textbf{Performance:} Escolha ponderada e otimizada dos tipos de dados (por exemplo, utilização do tipo \texttt{DECIMAL} para preços e valores monetários para evitar erros de arredondamento, e \texttt{VARCHAR} dimensionados adequadamente para texto).
\end{itemize}

\section{Entidades Principais}
O modelo de dados baseia-se num conjunto de entidades bem estruturadas que modelam a lógica de negócio do Cineteatro. A tabela seguinte apresenta a listagem das entidades principais e a respetiva descrição funcional das suas responsabilidades no sistema:

\begin{table}[h!]
\centering
\caption{Entidades do Sistema e Respetivas Responsabilidades}
\label{tab:entidades}
\vspace{0.2cm}
\renewcommand{\arraystretch}{1.3} % Melhora o espaçamento interno das linhas da tabela
\small
\begin{tabular}{ll}
\toprule
\textbf{\color{primaryColor}Entidade} & \textbf{\color{primaryColor}Responsabilidade Operacional} \\
\midrule
Pessoa      & Repositório de dados pessoais base partilhados (NIF, nome, contacto, morada). \\
Cliente     & Especialização de Pessoa; armazena o histórico e perfil de consumo. \\
Funcionário & Especialização de Pessoa; inclui dados de gestão de RH (função, salário). \\
Evento      & Registo completo de espetáculos com suporte a artistas múltiplos. \\
Bilhete     & Título de acesso emitido para um evento através de um operador da bilheteira. \\
Compra      & Registo de transações e consumo de produtos efetuadas nas áreas de bar. \\
Bar         & Identificação física do ponto de venda e controlo de inventário/stock. \\
\bottomrule
\end{tabular}
\end{table}

\paragraph{Notas sobre o modelo}
As entidades foram desenhadas para minimizar redundâncias e garantir que a informação pessoal é partilhada pela entidade `Pessoa` entre `Cliente` e `Funcionário`. As associações principais (FK) garantem rastreabilidade das transações ao cliente/funcionário responsável.

\section{Modelação de Dados}

\subsection{Diagrama Entidade-Relacionamento (DER)}
\begin{figure}[h!]
    \centering
    \includegraphics[width=0.85\textwidth]{DER_Cineteatro_v3.png}
    \caption{Diagrama Entidade-Relacionamento (DER) do Cineteatro (ficheiro: Images/DER_Cineteatro_v3.png — inserir no directório `Images/`)}
    \label{fig:der}
\end{figure}

\newpage

\subsection{Diagrama de Classes / Modelo Relacional (ER)}
\begin{figure}[h!]
    \centering
    \includegraphics[width=0.85\textwidth]{ER_Cineteatro_v3.png}
    \caption{Diagrama de Esquema Relacional (ER) do Cineteatro (ficheiro: Images/ER_Cineteatro_v3.png — inserir no directório `Images/`)}
    \label{fig:er}
\end{figure}

\section{SQL}

Os scripts desenvolvidos para a base de dados podem ser consultados diretamente:

\begin{itemize}
    \item \href{run:SQL/SQL_Cineteatro.sql}{SQL\_Cineteatro.sql}
    \item \href{run:SQL/insert_cineteatro.sql}{insert\_cineteatro.sql}
    \item \href{run:SQL/procedures_cineteatro.sql}{procedures\_cineteatro.sql}
    \item \href{run:SQL/functions_cineteatro.sql}{functions\_cineteatro.sql}
    \item \href{run:SQL/triggers.sql}{triggers.sql}
\end{itemize}

\newpage

\section{Forms}


\subsection{Bilheteira}

\begin{figure}[h!]
    \centering
    % imagem de interface da bilheteira (não encontrada no repositório)
    % \includegraphics[width=0.85\textwidth]{Images/PrintsDeBD1.png}
    \caption{Interface da Bilheteira (inserir print no ficheiro Images/PrintsDeBD1.png).}
    \label{fig:bilheteira}
\end{figure}

Nesta secção, é possível vender e remover bilhetes, bem como gerir os eventos associados ao sistema. A implementação baseia-se numa arquitetura desacoplada que utiliza \textit{Stored Procedures} para garantir a integridade dos dados.

\subsubsection{Gestão de Navegação e Dados}
O formulário principal utiliza um controlo de estado (\texttt{modoAtual}) para alternar dinamicamente a interface entre a visualização de bilhetes e eventos:

\begin{lstlisting}[language={[Sharp]C}, caption=Gestão de estados (excerto)]
// excerto: alterna modo e actualiza vista
ExibirBilhetes();
modoAtual = "Bilhetes";
// actualizar labels/botoes conforme o modo ativo
\end{lstlisting}

\subsubsection{Integridade e Remoção de Registos}
A remoção de dados é protegida por um tratamento de exceções robusto, capaz de capturar erros de integridade referencial (FK) do SQL Server:

\begin{lstlisting}[language={[Sharp]C}, caption=Tratamento de exceções e integridade referencial]
catch (SqlException ex)
{
    if (ex.Number == 547) // Erro de FK
    {
        MessageBox.Show("Nao e possivel remover este registo...");
    }
    else
    {
        MessageBox.Show(ex.Message);
    }
}
\end{lstlisting}

\subsubsection{Processo de Venda de Bilhetes}
A lógica de venda é encapsulada na \textit{Stored Procedure} \texttt{sp\_ComprarBilhete}. O formulário realiza primeiro validações locais (campos vazios ou tipos incorretos) antes de submeter o pedido à base de dados:

\begin{lstlisting}[language={[Sharp]C}, caption=Venda de bilhete (excerto)]
// validacoes: evento seleccionado, lugar e recibo válidos
cmd.CommandType = CommandType.StoredProcedure;
cmd.Parameters.AddWithValue("@Lugar", LugarInput.Text.Trim());
cmd.Parameters.AddWithValue("@ID_Evento", (long)EventoInput.SelectedValue);
// outros parametros (ID_Bilheteira, Data, ID_Recibo)
cmd.ExecuteNonQuery();
\end{lstlisting}


\newpage

\subsection{Recursos Humanos}
Esta secção resume a gestão dos funcionários no Cineteatro. O formulário de RH permite consultar a equipa, inserir novos colaboradores e remover registos com controlo de integridade, garantindo que os dados pessoais e profissionais permanecem consistentes.

\subsubsection{Funções principais}
Os métodos mais importantes são \texttt{CarregarFuncionarios()}, que mostra a equipa na grelha, \texttt{btnAdicionar_Click()}, que abre o formulário de inserção, e \texttt{btnRemover_Click()}, que elimina um funcionário apenas quando a base de dados o permite. O formulário \texttt{AdicionarFuncionario} garante ainda que a criação da pessoa e do funcionário é feita de forma transacional.

\begin{lstlisting}[language={[Sharp]C}, caption=Inserção de funcionário (excerto)]
// validacoes: campos obrigatorios; NIF e salário com formatos válidos
cmd.CommandType = CommandType.StoredProcedure;
cmd.Parameters.AddWithValue("@NIF", nifValor);
cmd.Parameters.AddWithValue("@Nome", nome.Text.Trim());
cmd.Parameters.AddWithValue("@Funcao", funcao.Text.Trim());
cmd.ExecuteNonQuery();
\end{lstlisting}

\newpage

\subsection{Bar}
O módulo do Bar é responsável por acompanhar o estado dos balcões, associar funcionários às áreas de venda e registar as transações de produtos. A interface foi desenhada para simplificar a operação diária e manter um histórico atualizado de consumo.

\subsubsection{Funções principais}
No Bar, \texttt{CarregarBares()} apresenta os registos existentes, \texttt{btnEditar_Click()} permite alterar o estado e o funcionário associado, e \texttt{btnCompras_Click()} abre o formulário \texttt{Venda}. Nesse formulário, \texttt{Venda_Load()} prepara os dados e \texttt{btnVender_Click()} grava a compra através da \textit{Stored Procedure} \texttt{sp\_RegistarCompraBar}.


\subsubsection{Venda no Bar}
A venda de produtos no bar é tratada pelo formulário \texttt{Venda}, que primeiro carrega a lista de produtos, configura os métodos de pagamento e mostra o histórico de transações. Antes de gravar a operação, o código valida o produto, a quantidade e o NIF do cliente, quando este é fornecido.

\begin{lstlisting}[language={[Sharp]C}, caption=Registo de venda no bar (excerto)]
// preparar parâmetros e executar procedimento que regista compra
cmd.CommandType = CommandType.StoredProcedure;
cmd.Parameters.AddWithValue("@ID_Bar", idBar);
cmd.Parameters.AddWithValue("@CodigoProduto", producto_select.SelectedValue);
cmd.Parameters.AddWithValue("@Quantidade", (int)numericUpDown1.Value);
cmd.ExecuteNonQuery();
\end{lstlisting}

\newpage
\subsection{Eventos}
O registo de eventos é um dos pontos mais importantes do sistema, porque combina validações locais com regras garantidas no SQL Server. O formulário liga-se à tabela de edifícios e envia os dados para a \textit{Stored Procedure} que insere o evento.

\begin{lstlisting}[language={[Sharp]C}, caption=Registo de eventos com validação local]
private void btnGravar_Click(object sender, EventArgs e)
{
    if (string.IsNullOrWhiteSpace(NomeInput.Text))
        return;

    if (!TimeSpan.TryParse(DuracaoInput.Text, out TimeSpan duracao))
        return;

    using (SqlCommand cmd = new SqlCommand("dbo.sp_AdicionarEvento", conn))
    {
        cmd.CommandType = CommandType.StoredProcedure;
        cmd.Parameters.AddWithValue("@Nome", NomeInput.Text.Trim());
        ...
        cmd.ExecuteNonQuery();
    }
}
\end{lstlisting}

\begin{lstlisting}[language={[Sharp]C}, caption=Registo de evento (excerto)]
// validações locais: Nome e duração válidos
cmd.CommandType = CommandType.StoredProcedure;
cmd.Parameters.AddWithValue("@Nome", NomeInput.Text.Trim());
cmd.Parameters.AddWithValue("@Duracao", duracao);
cmd.ExecuteNonQuery();
\end{lstlisting}


\section{Conclusão}
O sistema implementado satisfaz os requisitos funcionais principais e adopta boas práticas de modelação e integridade referencial. A utilização de stored procedures centraliza as regras críticas no servidor, simplificando a lógica de cliente e facilitando auditoria.



\end{document}
